% !TEX root = ../main.tex
\chapter{大面积单模设计原则}
\label{ch:single-mode-design}

\section*{学习目标}
\begin{itemize}
  \item 理解为什么“大面积”和“单模”在 PCSEL 中天然存在张力，以及这种张力为什么不能只靠单胞带结构解决。
  \item 掌握模鉴别必须同时发生在单胞、有限孔阵和器件工作点三个层次上的原因。
  \item 明确哪些设计问题主要由光学仿真回答，哪些问题必须进入电学与热学模型后才能判断。
\end{itemize}

\section*{先修知识}
建议已掌握 \cref{ch:rcwa,ch:fdtd-fem,ch:epitaxial-optics,ch:electrical,ch:eto-coupling} 中关于开放结构、有限尺寸、外延层、电注入和热反馈的内容。

\section*{本章路线图}
\ChapterModelLevel{L2→L4}
前面三章已经把一个关键事实说清楚了：真实 PCSEL 的目标并不是在被动冷腔里找到一个看起来漂亮的 Gamma 点模式，而是在给定电流、温度和外延约束下，让一个目标模式在大面积孔阵中稳定地赢过所有竞争模式。这个目标决定了本章的写法。我们先解释为什么大面积单模从来不是“把面积做大而模态不变”这么简单，再建立模鉴别裕量的概念，随后把设计原则分成三个尺度来讨论：单胞尺度的辐射与对称性、有限孔阵尺度的横向包络与边缘损耗、器件尺度的注入与热分布。最后再把这些原则翻译成真正可执行的仿真任务，并说明在 Lumerical 与 COMSOL 风格环境中通常各自承担哪些计算。

\section{为什么大面积单模是 PCSEL 最难也最值钱的问题}
PCSEL 之所以受到持续关注，一个核心原因就是它有潜力同时获得较大出光面积和较好光束质量。这个说法的潜台词是：目标模式必须在较大横向面积上保持相干，而不是把许多局域振荡区域松散地拼在一起。问题就在这里。器件面积一旦增大，可利用的横向自由度也随之增多，竞争模数量、边缘效应、注入不均匀性和热非均匀性都会变得更显著。于是“大面积”不是一个被动的几何放大，而是主动增加了模式竞争的维数。
实现大面积相干的关键在于建立有效的横向耦合物理机制。\cref{fig:ch20_array_coupling}总结了三种主要的同步途径：倏逝波隧穿耦合（短程最近邻作用）、辐射场耦合（长程远场干涉）以及共享全局腔模（集体约束条件）。实际设计中往往需要综合运用这三种机制，通过适度密集的排布增强倏逝耦合，同时在边缘设置分布式反馈结构提供额外的模式选择性。

% !TEX root = ../main.tex
% Figure: Large-area coherent array synchronization mechanisms
% Chapter: ch20 - Large-Area Single-Mode Operation Strategies
% Description: Different coupling pathways for phase locking

\begin{figure}[htbp]
  \centering
  \begin{tikzpicture}[scale=0.92]
    % Array of emitters (simplified as circles)
    \def\emitterradius{0.3}
    
    % 5x5 emitter grid
    \foreach \x in {-4,-2,...,4} {
      \foreach \y in {-4,-2,...,4} {
        % Emitting aperture
        \fill[white] (\x,\y) circle (\emitterradius);
        \draw[thick, blue] (\x,\y) circle (\emitterradius);
        
        % Phase indicator (arrow showing oscillation phase)
        \pgfmathrandominteger{\phase}{0}{359}
        \draw[->, red, thick] (\x,\y) -- ++(\phase:0.2);
      }
    }
    
    % Background: Common substrate
    \fill[gray!10] (-5.5,-5.5) rectangle (5.5,5.5);
    \node[above right, gray] at (-5.5,-5.5) {公共衬底};
    
    % Coupling mechanism 1: Evanescent tunneling (nearest neighbors)
    \draw[very thick, green!60!black, <->] (-2,0) -- node[above] {\scriptsize 倏逝耦合} (0,0);
    \draw[very thick, green!60!black, <->] (0,0) -- (2,0);
    \draw[very thick, green!60!black, <->] (0,-2) -- (0,0);
    \draw[very thick, green!60!black, <->] (0,0) -- (0,2);
    
    % Coupling mechanism 2: Radiation coupling (long-range)
    \draw[very thick, orange, ->, bend left=20] (-4,-4) to[out=30,in=150] (4,4);
    \draw[very thick, orange, ->, bend right=20] (4,-4) to[out=150,in=30] (-4,4);
    \node[orange, align=center] at (0,3.5) {辐射耦合\\远场重叠};
    
    % Coupling mechanism 3: Shared cavity mode (global)
    \draw[ultra thick, purple, dashed] (0,0) circle (5);
    \node[purple, above] at (0,5.2) {共享腔模};
    
    % Boundary conditions
    \draw[double distance=2pt, thick] (-5.5,-5.5) rectangle (5.5,5.5);
    \node[below, font=\small] at (0,-5.8) {金属包覆边界或 DBR反射镜};
    
    % Zoom-in inset: Unit cell detail
    \begin{scope}[shift={(7,3)}, scale=2]
      % Single emitter with PhC cladding
      \fill[white] (0,0) circle (0.3);
      \draw[thick, blue] (0,0) circle (0.3);
      
      % Surrounding photonic crystal holes
      \foreach \angle in {0,60,...,300} {
        \fill[gray] (\angle:0.6) circle (0.08);
      }
      
      % Fundamental mode profile
      \draw[very thick, red, domain=0:360, samples=100] 
        plot({0.2*cos(\x)*(1 + 0.12*cos(2*\x))},{0.2*sin(\x)*(1 + 0.12*cos(2*\x))});
      
      \node[above, font=\small] at (0,0.8) {单胞放大};
      \node[below, align=center, font=\footnotesize] at (0,-0.8) {基模$\mathrm{LP}_{01}$\\近场高斯型};
    \end{scope}
    
    % Coherence area illustration
    \fill[yellow!20, opacity=0.5] (-1,-1) circle (1.5);
    \node[yellow!60!black, align=center] at (-1,-1) {相干区域\\$A_c$};
    
    % Key parameters box
    \node[draw, align=left, font=\small, anchor=north west] at (-6,6) {
      \textbf{同步条件}:\\
      $\bullet$ 耦合强度$\kappa >$频率失谐$\Delta\omega$\\
      $\bullet$ 相位锁定时间$\tau_c < $光子寿命$\tau_p$\\
      $\bullet$ 阵列尺寸$L <$相干长度$L_c$
    };
    
    % Output beam quality metrics
    \node[draw, align=left, font=\small, anchor=south east] at (6,-6) {
      \textbf{光束质量指标}:\\
      $\bullet$ $M^2 < 1.2$接近衍射极限\\
      $\bullet$ 旁瓣抑制比$>15\,\mathrm{dB}$\\
      $\bullet$ 指向稳定性$<0.1^\circ$
    };
  \end{tikzpicture}
  \caption{大面积 PCSEL 阵列中相干锁定的三种主要横向耦合机制：倏逝近邻耦合、辐射远场耦合以及共享全局腔模。右上角给出单胞局域模式示意，右下角列出评价大面积单模器件时常用的光束质量指标。}
  \label{fig:ch20_array_coupling}
\end{figure}


如果只从二维光子晶体能带的角度看问题，很容易以为只要选中了 Gamma 点附近一支低损耗模，单模问题就已经基本解决。事实上，这最多只解决了第一道门槛：说明在无限周期、冷腔、小信号框架下，目标模有希望成为候选模。真正的器件级单模要求远高于此。它要求目标模不仅阈值低，还要对有限孔阵边界、局部增益起伏、温度漂移和偏振扰动具有足够鲁棒性。换句话说，单模不只是“首先起振”，而是“在工作窗口内持续压住其它模”。

\begin{IntuitionBox}
对大面积 PCSEL 来说，最危险的误解是把“存在一个低损耗目标模”误认为“器件会稳定单模工作”。前者只是候选资格，后者则是一个跨越光学、电学和热学的稳定性结论。
\end{IntuitionBox}

\section{模鉴别不是一句话，而是一个裕量问题}
为了让后面的讨论有共同语言，可以把目标模 $m_\star$ 相对于竞争模 $m_j$ 的优势写成净模裕量
\begin{equation}
\Delta G_{\star j}(I,T)
=
\bigl[g_{\mathrm{mod},\star}(I,T)-\alpha_{\star}(I,T)\bigr]
-
\bigl[g_{\mathrm{mod},j}(I,T)-\alpha_{j}(I,T)\bigr],
\label{eq:modal_margin}
\end{equation}
其中 $g_{\mathrm{mod},i}$ 是第 $i$ 个模在给定工作点下的模增益，$\alpha_i$ 是其总损耗，包含辐射损耗、内部吸收、边缘损耗以及工作状态下可能出现的附加损耗。\cref{eq:modal_margin} 不是新的基本定律，它是一个\textbf{组织设计问题的记账式写法}。它的价值在于提醒我们：

第一，单模判据是相对的，不是绝对的。目标模本身阈值低还不够，还必须比竞争模更有优势。

第二，这个优势依赖工作点 $(I,T)$。也就是说，单模不是结构固有的单一标签，而是工作窗口中的一个稳定区间结论。

第三，\cref{eq:modal_margin} 中的每一项可能来自不同物理模块。$g_{\mathrm{mod},i}$ 需要量子阱增益与注入分布，$\alpha_i$ 则既依赖开放光学，也依赖掺杂吸收、金属邻近损耗、边缘损耗和热漂移。

因此，大面积单模设计的实质不是盯住某一个高 $Q$ 数字，而是尽可能在目标工作窗口内让
\begin{equation}
\Delta G_{\star j}(I,T) > 0,
\qquad \forall j\neq \star,
\label{eq:single_mode_margin}
\end{equation}
并且这个不等式对工艺误差和热漂移不至于过分敏感。

\begin{RigorousBox}
\cref{eq:modal_margin,eq:single_mode_margin} 属于\textbf{器件级组织框架}。它不是仅靠冷腔本征解就能完全求出的量，因为其中至少一半信息来自注入分布与温升。用被动 $Q$ 或被动线宽去直接替代 $\Delta G_{\star j}$，属于层级偷换。
\end{RigorousBox}

\begin{table}[htbp]
  \centering
  \footnotesize
  \caption{模裕量各项与求解模块的交接关系。}
  \label{tab:modal_margin_handoff}
  \begin{tabularx}{\textwidth}{@{}p{2.6cm}p{3.0cm}Y@{}}
    \toprule
    裕量中的量 & 主要来源层级 & 更自然的求解接口 \\ \midrule
    $g_{\mathrm{mod},i}$ & 有源区、注入分布、温度回写后的模增益 & 量子阱增益模型 + 电学/热学回写 + 模场重叠 \\
    $\alpha_{\mathrm{rad},i}$ & 开放辐射损耗 & RCWA，或有限阵列 FDTD/FEM \\
    $\alpha_{\mathrm{int},i}$ & 自由载流子吸收、金属损耗、材料吸收 & 外延层与器件损耗模型 \\
    $\alpha_{\mathrm{edge},i}$ & 有限孔阵边缘泄漏与侧向逃逸 & 有限阵列全波，或带包络的有限尺寸 CWT \\
    热回写项 & $\Delta T\rightarrow\Delta n,\Delta g,\Delta\alpha$ & 热-电-光耦合模块 \\
    \bottomrule
  \end{tabularx}
\end{table}

\section{第一尺度：单胞与 Gamma 点物理决定的是候选资格}
设计的第一尺度是单胞尺度。这里关心的问题是：在无限周期近似下，哪类 Gamma 点模具有较低被动辐射损耗，哪类近场对称性更容易给出期望偏振和远场，哪类对称性破缺或双晶格扰动能把目标模与竞争模分开。这一层上的典型设计变量包括晶格常数 $a$、孔半径 $r$、占空比、刻蚀深度、双晶格位移、单胞形状微扰和上下层折射率对比等。

这一层可以回答的，是“目标模是否值得继续研究”。如果在单胞尺度上目标模就已经没有明显辐射优势，或者偏振选择完全不稳定，那么继续投入有限器件和电热建模通常是浪费时间。但必须认识到，单胞尺度只给出候选资格，不给出最终单模结论。因为单胞尺度默认了无限周期、相同局域环境和均匀增益代理，这些假设一旦进入大面积器件就都会受到挑战。

从求解器角度看，这一层最自然的工具是 \cref{ch:pwe,ch:rcwa} 中的 PWE、CWT 和 RCWA。如果用商业软件来承载这一层，Lumerical 中通常优先使用 MODE/FDE 做纵向等效折射率与波导模检查，再用 RCWA 或周期 FDTD 做单胞散射与辐射通道分析；COMSOL 中通常使用 Wave Optics 的频域周期单胞模型，通过 Floquet 型相位周期边界和端口或入射波设置读取反射、透射、衍射效率与局域场分布。

\section{第二尺度：有限孔阵与边缘损耗决定的是横向可实现性}
一旦走出无限周期近似，问题立刻变了。真实 PCSEL 只有有限个周期，边界会引入额外泄漏，目标模也会从严格 Bloch 模转化为带有横向包络的有限腔模式。此时设计的第二尺度就出现了：横向包络如何被组织，边缘处是增强还是抑制目标模，有限孔阵尺寸是否已经足以支撑目标模的相位一致性。

所谓边缘损耗，本质上就是原本在无限周期里被周期散射网络反复折返的能量，到了有限边界处不再被完整送回，而是沿侧向或斜向辐射通道逃逸出去。因此它不是“附加罚项”，而是 Bragg 反馈网络被截断后的直接物理后果。

这一层最重要的直觉是，大面积并不自动有利于单模。面积增大确实可以减小衍射发散、提高功率潜力，但同时也会给许多横向模提供生存空间。为了让目标模在大面积下仍然占优，常见做法不是简单把阵列做大，而是同时设计横向损耗分布。例如：
\begin{itemize}
  \item 通过边缘区扰动增加竞争模的横向泄漏；
  \item 通过双晶格或对称性工程让目标模在近法向辐射上更占优势；
  \item 通过局部参数渐变或包络控制，让目标模横向分布与注入区域更匹配；
  \item 在保持单模性的前提下控制上下辐射和侧向泄漏的权衡。
\end{itemize}

这里的关键不是“损耗越小越好”，而是“目标模与竞争模的损耗排序是否朝正确方向移动”。有时为了压制竞争模，反而会故意让整体损耗略有增加，只要目标模的相对优势扩大，器件级单模性就可能改善。

这一层通常已经超出周期单胞 RCWA 的默认边界，需要转向有限尺寸全波仿真。Lumerical 中最常见的是 FDTD；COMSOL 中则更常见的是 Wave Optics 的三维频域或本征频率求解并辅以 PML。这里的输出应当包括有限阵列的近场、远场、复频率或等效 $Q$、边缘敏感性，以及目标模与竞争模的排序是否随阵列大小改变。

\section{第三尺度：注入与温升决定的是工作点下的真实胜负}
即使前两个尺度都看起来很好，器件仍然可能在实际工作中丢掉单模性。原因在于第三尺度，也就是器件工作尺度。到了这里，模式竞争不再仅仅由被动辐射损耗决定，而是由载流子分布、局部模增益、串联电阻、Joule 热、非辐射复合热以及温升引起的折射率和增益谱漂移共同决定。

这正是 \cref{ch:electrical,ch:eto-coupling} 反复强调的地方。目标模光学上“形状好看”，并不意味着它在电流注入下得到的净增益最大；目标模冷腔频率“居中”，也不意味着温升后它仍然处在增益谱最有利的位置。大面积器件尤其容易出现这类反转，因为横向注入不均匀和热非均匀会随面积扩大而加剧，导致目标模原本依赖的对称性被工作状态主动破坏。

因此，真正的大面积单模设计必须把外延层光学、注入路径和热扩散一起作为设计变量。常见的协同思路包括：
\begin{itemize}
  \item 让量子阱重叠最大的区域尽量与电流高密度区域对齐，而不是错位；
  \item 通过电极形状和限制层设计，减少边缘过强注入导致的竞争模局部起振；
  \item 通过层栈和散热路径设计，降低温升诱发的模排序翻转；
  \item 让目标模对小幅折射率漂移和增益峰位漂移具有更好的鲁棒性。
\end{itemize}

这部分若采用 Lumerical 工作流，通常需要 DEVICE 中的 CHARGE / HEAT 与光学求解器之间做脚本化数据交换；若采用 COMSOL 工作流，则更常见的是在 Wave Optics、Semiconductor 与 Heat Transfer 模块之间做统一几何下的多物理场耦合或分阶段迭代。哪种方式更优不是定理，而是工程权衡：Lumerical 的光学链往往更高效，COMSOL 的自定义耦合 PDE 链往往更灵活。

\section{边缘损耗与表面提取不是对立词，而是同一预算的两端}
很多新手在谈单模设计时，会把“提高表面提取效率”和“抑制边缘泄漏”看成两个互不相干的目标。实际上，这二者共享同一份辐射预算。PCSEL 需要垂直辐射才有面发射输出，但又不希望这种辐射机制让竞争模也轻易起振；同时，器件还不希望能量在横向边缘无序泄漏，因为那会破坏单模相干面积。于是设计问题不再是简单追求最小辐射损耗，而是把有限的辐射耦合尽可能分配给目标模需要的垂直输出通道，而不是分配给竞争模或边缘泄漏通道。

这就是为什么一些看起来“很高 Q”的结构在大面积器件里并不一定表现最好。若过分压低所有辐射损耗，虽然被动 $Q$ 可能很高，但目标模未必最容易提取，竞争模也可能被一并保护起来。相反，适度引入受控垂直辐射并选择性增加竞争模边缘泄漏，反而更有利于稳定单模工作。

\begin{ExampleBox}
设两种方案 A 与 B 在单胞层面都能支持目标 Gamma 点模。方案 A 的被动 $Q$ 更高，但有限阵列中竞争模的边缘泄漏也同样很小；方案 B 的目标模被动 $Q$ 稍低，却能通过边缘扰动显著增大竞争模的侧向泄漏。在器件级上，方案 B 反而可能更容易保持大面积单模，因为它扩大的是目标模相对于竞争模的净裕量，而不是盲目提高所有模的寿命。
\end{ExampleBox}

\section{把设计原则翻译成仿真任务：应该在哪类工具里算什么}
到了这里，设计原则必须落回可计算任务。最实用的做法，是按“物理问题”而不是“软件名称”来分配求解器，然后再把这些求解器映射到 Lumerical 或 COMSOL 中。\cref{tab:singlemode-toolmap} 给出的是一种实用的、但并非唯一的映射方式。

\begin{table}[htbp]
  \centering
  \footnotesize
  \caption{大面积单模设计中常见问题与 Lumerical / COMSOL 的典型实现映射。这里给出的是工程上常见分工，而不是唯一正确方案。}
  \label{tab:singlemode-toolmap}
  \setlength{\tabcolsep}{4pt}
  \begin{tabularx}{\textwidth}{@{}p{2.2cm}p{1.8cm}Y Y@{}}
    \toprule
    待判断性质 & 首选物理层级 & Lumerical 中常见实现 & COMSOL 中常见实现 \\
    \midrule
    纵向波导模、$n_{\mathrm{eff}}$、约束因子 & 2D / 2.5D 光学 & MODE/FDE 计算层栈模场、有效折射率与损耗 & 2D 横截面 Wave Optics 本征/模态分析 \\
    单胞辐射、偏振、上下出光比例 & 周期开放光学 & RCWA 或周期 FDTD & Wave Optics 频域单胞，配合相位周期边界与端口/入射波 \\
    有限孔阵模场、远场、边缘敏感性 & 有限尺寸全波 & FDTD，配合 PML 与近远场变换 & 3D Wave Optics 频域或本征频率研究，配合 PML \\
    电流扩展、载流子分布、J-V & 漂移-扩散电学 & CHARGE & Semiconductor \\
    温升、热源、热漂移 & 热传导 & HEAT & Heat Transfer in Solids \\
    热-电-光迭代后工作窗口 & 多物理场耦合 & 脚本串联 RCWA/FDTD 与 CHARGE/HEAT & Wave Optics、Semiconductor 与 Heat Transfer 分阶段耦合或统一模型 \\
    \bottomrule
  \end{tabularx}
\end{table}

如果进一步问“简单建模过程是什么”，可以先给出最小可执行版本：

在 Lumerical 路径中，先用 MODE/FDE 检查纵向模与重叠，再用 RCWA 扫单胞参数和法向附近角谱，得到候选结构；随后把候选结构放入 FDTD 的有限阵列模型里计算近场、远场和边缘敏感性；最后在 CHARGE 中解注入分布、在 HEAT 中解温升，并把得到的 $N(\mathbf{r})$ 与 $T(\mathbf{r})$ 以代理方式回写到光学模型中做阈值和模竞争再评估。

在 COMSOL 路径中，先在二维横截面或简化三维模型中做纵向模分析；随后在单胞三维频域模型中用相位周期边界和端口求出反射、透射与辐射分配；再切换到有限阵列三维模型，用 PML 包围开放区域求近场和远场；最后在同一几何下加入 Semiconductor 与 Heat Transfer 物理接口，解出电流密度、载流子密度和温度场，再将这些结果通过折射率和增益代理回写到光学接口中迭代。

\section{真正的设计目标是工作窗口，而不是最佳单点}
教材最后必须把一句话说死：大面积单模设计追求的不是一个“最优参数点”，而是一个可工作的参数窗口。若某个结构只有在极窄的孔半径、公差和电流点上保持单模，它在研究原型中可能还能看到漂亮图样，但在真实器件中通常缺乏可复现性。相反，哪怕某个方案的理论最优值略差，只要它对工艺偏差、热漂移和注入不均匀有更强鲁棒性，它就更接近真正可用的设计。

这也是为什么本章始终围绕 \cref{eq:modal_margin} 展开，而不是围绕某个单独的 $Q$、某个单独的远场主瓣角或某个单独的阈值增益。对大面积 PCSEL 来说，单模性本质上是一个多约束交集问题，只有把这些约束压进同一工作窗口，结果才有器件意义。

\begin{PitfallBox}
“把被动 $Q$ 最高的方案选出来，再交给电学和热学去补救”通常不是高效流程。真正高效的做法，是在单胞和有限阵列阶段就开始追踪目标模相对于竞争模的排序，而不是等器件级耦合完成后才发现一开始选错了候选结构。
\end{PitfallBox}

\section*{本章小结}
大面积单模设计的核心，不是追求一个孤立的低损耗目标模，而是在单胞、有限孔阵和器件工作点三个尺度上同时扩大目标模相对于竞争模的净裕量。单胞尺度决定候选资格，有限阵列尺度决定横向可实现性，器件尺度决定真实工作点下的胜负。Lumerical 与 COMSOL 都可以承载这条链，但各自更擅长的部分不同：前者通常更高效地承担 RCWA/FDTD/模式分析，后者通常更灵活地承担多物理场 PDE 耦合。无论使用哪条路径，最终目标都不是最佳单点，而是稳定、可复现的工作窗口。

\section*{练习题}
\begin{enumerate}
  \item \basicq 结合 \cref{eq:modal_margin} 解释为什么“目标模阈值低”并不自动等价于“器件单模稳定”。

  \item \basicq 说明为什么大面积器件中的边缘工程应当被视为模鉴别设计，而不是几何收尾工作。

  \item \midq 选择一种你熟悉的结构扰动方式，分析它更可能改变的是单胞辐射排序、有限阵列边缘损耗，还是电热工作点下的净模增益排序。

  \item \hardq 分别写出一条以 Lumerical 为主和一条以 COMSOL 为主的大面积单模验证路径，并指出各自最关键的接口变量是什么。
\end{enumerate}

\section*{建议继续阅读}
\begin{itemize}
  \item 下一章将进一步讨论名义值设计为何仍然不够，并系统引入有限尺寸、无序与工艺误差鲁棒性的问题，把“单模设计原则”升级成“可制造器件设计原则”。 这条阅读的重点是把本章的输出顺着主线接到后续模块，而不是停成孤立知识点。

  \item 若想继续深化模式竞争与单模稳定性，可结合 \cref{ch:robustness,ch:workflow} 一起阅读，把本章的设计原则放入完整验证与容差分析工作流中。 这条阅读的重点是把本章的输出顺着主线接到后续模块，而不是停成孤立知识点。
\end{itemize}


